\newpage
\begin{center}
	\large\textbf{Минобрнауки России}
	
	\large\textbf{Юго-Западный государственный университет}
	\vskip 1em
	\normalsize{Кафедра программной инженерии}
	\vskip 1em
	\ifВКР{
		\begin{flushright}
			\begin{tabular}{p{.4\textwidth}}
				\centering УТВЕРЖДАЮ: \\
				\centering Заведующий кафедрой \\
				\hrulefill \\
				\setarstrut{\footnotesize}
				\centering\footnotesize{(подпись, инициалы, фамилия)}\\
				\restorearstrut
				«\underline{\hspace{1cm}}»
				\underline{\hspace{3cm}}
				20\underline{\hspace{1cm}} г.\\
			\end{tabular}
		\end{flushright}
	}\fi
\end{center}
\vspace{1em}
\begin{center}
	\large
	\ifВКР{
		ЗАДАНИЕ НА ВЫПУСКНУЮ КВАЛИФИКАЦИОННУЮ РАБОТУ
		ПО ПРОГРАММЕ БАКАЛАВРИАТА}
	\else
	ЗАДАНИЕ НА КУРСОВУЮ РАБОТУ (ПРОЕКТ)
	\fi
	\normalsize
\end{center}
\vspace{1em}
{\parindent0pt
	Студента \АвторРод, шифр\ \Шифр, группа \Группа
	
	1. Тема «\Тема\ \ТемаВтораяСтрока»
	\ifВКР{
		утверждена приказом ректора ЮЗГУ от \ДатаПриказа\ № \НомерПриказа
	}\fi.
	
	2. Срок предоставления работы к защите \СрокПредоставления
	
	3. Исходные данные для создания программной системы:
	
	3.1. Перечень решаемых задач:}

\renewcommand\labelenumi{\theenumi)}
\begin{enumerate}
	\item провести анализ предметной области йога-студий и существующих решений для записи на занятия;
	\item разработать концептуальную модель мобильного приложения «Йога-студия», определить основные сущности и связи между ними;
	\item спроектировать архитектуру приложения (Single Activity + Fragments) и пользовательский интерфейс;
	\item разработать и протестировать мобильное приложение для платформы Android на языке Kotlin с использованием SharedPreferences.
\end{enumerate}

{\parindent0pt
	3.2. Входные данные и требуемые результаты для программы:}

\begin{enumerate}
	\item \textbf{Входными данными} для программной системы являются: регистрационные данные пользователей (имя, email, пароль), профильные данные пользователей (аватар, город, уровень подготовки), данные справочников занятий и инструкторов, параметры записи на занятия, текстовые сообщения отзывов и оценки, поисковые запросы и параметры фильтрации по городу.
	
	\item \textbf{Выходными данными} для программной системы являются: сформированный каталог занятий с фильтрацией по городу, расписание с календарём, профиль пользователя с аватаром и настройками, список записей пользователя, отзывы и рейтинг занятий, статистика прогресса в виде графиков, уведомления о подтверждении записи и напоминания о занятиях.
\end{enumerate}

{\parindent0pt
	
	4. Содержание работы (по разделам):
	
	4.1. Введение.
	
	4.2. Анализ предметной области.
	
	4.3. Техническое задание: основание для разработки, назначение разработки, требования к программной системе, требования к оформлению документации.
	
	4.4. Технический проект: общие сведения о программной системе, проект данных программной системы, проектирование архитектуры программной системы, проектирование пользовательского интерфейса программной системы.
	
	4.5. Рабочий проект: спецификация компонентов и классов программной системы, тестирование программной системы.
	
	4.6. Заключение.
	
	4.7. Список использованных источников.
	
	4.8. Приложения (листинги программного кода, плакаты).
	
	5. Перечень графического материала:
	
	\списокПлакатов
	
	\vskip 2em
	\begin{tabular}{p{6.8cm}C{3.8cm}C{4.8cm}}
		Руководитель \ifВКР{ВКР}\else работы (проекта) \fi & \lhrulefill{\fill} & \fillcenter{Т. Н. Конаныхина}\\
		\setarstrut{\footnotesize}
		& \footnotesize{(подпись, дата)} & \footnotesize{(инициалы, фамилия)}\\
		\restorearstrut
		Задание принял к исполнению & \lhrulefill{\fill} & \fillcenter{В. Е. Бабичев}\\
		\setarstrut{\footnotesize}
		& \footnotesize{(подпись, дата)} & \footnotesize{(инициалы, фамилия)}\\
		\restorearstrut
	\end{tabular}
}

\renewcommand\labelenumi{\theenumi.}

\newpage